
\subsubsection{5.1 H-E1: All 5 Gates Pass}

\textbf{Table 1: FAIL FAST Gate Results (source: h-e1/experiment\_results.json)}

\begin{table*}[htbp]
\centering
\resizebox{\dimexpr 2\columnwidth+\columnsep\relax}{!}{%
\begin{tabular}{lllll}
\toprule
Gate & Metric & Value & Threshold & Status \\
\midrule
G0 & Join coverage & 0.862 & $\geq$ 0.80 & PASS \\
G1 & partial $r^{2}(log$\_count\_z) & 0.605 & > 0.01 & PASS \\
G2 & partial $r^{2}(diversity$\_ratio\_z) & 0.975 & > 0.01 & PASS \\
G3 & std(diversity\_ratio) & 0.246 & > 0.10 & PASS \\
G4 & max VIF & 2.14 & < 10.0 & PASS \\
\bottomrule
\end{tabular}
}
\end{table*}

Additional: Pearson r(log\_count\_z, diversity\_ratio\_z) = $-0.324$ (collinearity failsafe not triggered). 75 of 87 benchmarks matched in the fuzzy join; 67 of 87 have diversity statistics computed.

The FAIL FAST all-pass establishes: (a) the diversity signal exists and is joinable (G0 = 0.862); (b) both predictors are genuinely time-independent, not temporal proxies (G1 = 0.605, G2 = 0.975, both far exceeding the 0.01 threshold); (c) adequate cross-benchmark variance for Cox regression (G3 std = 0.246, 2.$5\times$ above threshold). Any downstream null is therefore attributable to genuine absence of effect rather than data quality failure.

\begin{figure}[htbp]
\centering
\includegraphics[width=\columnwidth]{figures/coverage_heatmap.png}
\caption{Coverage Heatmap}
\label{fig:coverage_heatmap}
\end{figure}

\textit{Figure 3: Coverage heatmap --- 75 of 87 benchmarks matched in the fuzzy join (Gate G0 = 0.862). Twelve benchmarks fell below token\_sort\_ratio = 85.}

\subsubsection{5.2 H-M1: Near-Perfect Null}

\textbf{Table 2: Cox Regression Results (source: h-m1/experiment\_results.json)}

\begin{table*}[htbp]
\centering
\resizebox{\dimexpr 2\columnwidth+\columnsep\relax}{!}{%
\begin{tabular}{llllll}
\toprule
Metric & Value & Threshold & Status &  &  \\
\midrule
LRT statistic & 0.0040 & --- & --- &  &  \\
LRT p-value & 0.9495 & < 0.05 & FAIL &  &  \\
Hazard Ratio (HR) & 1.006 & \ & $HR-1$\ & $\geq$ 0.10 & FAIL \\
95\% CI & [0.846, 1.196] & excludes 1.0 & FAIL &  &  \\
\ & $HR-1$\ &  & 0.0056 & $\geq$ 0.10 & FAIL (~$18\times$ below threshold) \\
M0 log-likelihood & $-750.7335$ & --- & --- &  &  \\
M1 log-likelihood & $-750.7315$ & --- & --- &  &  \\
$\Delta logL$ & 0.0020 & --- & --- &  &  \\
Concordance index (M1) & 0.7363 & --- & --- &  &  \\
Complete-case rows & 258 & --- & 87 NaN dropped &  &  \\
Displacement events & 176 & --- & 68.2\% event rate &  &  \\
\bottomrule
\end{tabular}
}
\end{table*}

HR = 1.006 is indistinguishable from no effect. $\Delta logL$ = 0.0020 --- the predictor adds essentially zero explanatory power to the baseline model. The 95\% CI [0.846, 1.196] straddles 1.0 by a wide margin. Both H1 (lock-in, requires HR < 1) and H2 (saturation, requires HR > 1) are falsified by the pre-specified success criteria. Direction: H0 (null).

\textbf{Model validity:} Concordance index = 0.7363 indicates the baseline covariates (task\_age, log\_publication\_volume) do predict displacement timing, confirming the model is functional. The diversity predictor is uninformative given these controls.

\begin{figure}[htbp]
\centering
\includegraphics[width=\columnwidth]{figures/forest_plot.png}
\caption{Forest Plot}
\label{fig:forest_plot}
\end{figure}

\textit{Figure 4: Forest plot showing HR = 1.006 (95\% CI = [0.846, 1.196]) for the diversity predictor. The confidence interval straddles 1.0 by a wide margin in both directions.}

\subsubsection{5.3 Kaplan-Meier Analysis}

KM curves for Q1 (lowest diversity) and Q4 (highest diversity) show no meaningful separation visually.

\begin{figure}[htbp]
\centering
\includegraphics[width=\columnwidth]{figures/km_quartiles.png}
\caption{Kaplan-Meier Q1 vs Q4}
\label{fig:km_quartiles}
\end{figure}

\textit{Figure 5: Kaplan-Meier survival curves for Q1 (lowest submission-count diversity) vs Q4 (highest submission-count diversity). No separation is apparent, consistent with the Cox null result. No numeric log-rank p-value was extracted in this pipeline run (see Limitation L6).}

\subsubsection{5.4 Relationship to Prior Directional Signal}

\textbf{Table 3: Comparison of Two Distinct Predictor Tests}

\begin{table}[htbp]
\centering
\resizebox{\columnwidth}{!}{%
\begin{tabular}{lll}
\toprule
Property & h-m1 Run 2 (prior, archived) & H-M1 (current) \\
\midrule
Predictor & $\Delta score$\_lag1\_z (*score velocity*) & log\_unique\_paper\_count\_at\_intro\_z (*submission count*) \\
Rows used & ~64 & 258 \\
Displacement events & ~22 & 176 \\
HR & 0.871 & 1.006 \\
95\% CI & (wide, not reported) & [0.846, 1.196] \\
LRT p & 0.565 & 0.9495 \\
Interpretation & Directional signal (underpowered) & Near-perfect null (adequate power) \\
\bottomrule
\end{tabular}
}
\end{table}

These two runs test \textit{different constructs}. The prior result (HR = 0.871) is for \textit{score velocity} (how fast SOTA scores improved); the current result is for \textit{submission-count diversity} (how many distinct papers submitted). The current null does not supersede or invalidate the prior directional signal for score velocity --- that predictor has not been tested at adequate statistical power. The two results are complementary: submission-count breadth is definitively null; score velocity remains an open question.

